\section{实验}
\label{sec:exp}

\subsection{实验设置}
\label{sec:exp:setup}

\paragraph{主干模型。}
除特别说明外，我们使用 \textsc{OpenVLA-7b}（\texttt{openvla/openvla-7b}）。\textsc{OpenVLA-OFT}（\texttt{moojink/openvla-7b-oft-\allowbreak finetuned-libero-spatial}）用作方案特有的负对照（\S\ref{sec:analysis:oft}）。外部对齐 LLM 为 \texttt{NousResearch/\allowbreak Llama-2-7b-chat-hf}，是与 Meta Llama-2-7b-chat 权重相同的非门控镜像。

\paragraph{评测套件。}
\textbf{LIBERO-Long}（10 个长时序操作任务，每任务 50 条轨迹），用作良性任务和视觉场景的来源；\textbf{SafeAgentBench} \cite{safeagentbench2024}，我们从中选取 120 个具身危害子集\footnote{使用 750 任务 AI2-THOR 分集以保证流水线兼容性。}；以及\textbf{指令条件化 AdvBench} 分集，将从我们策划分类体系（暴力、儿童安全、火/电、化学、财产损毁、疏忽；40 条提示 × 5 次改写）中提取的 128 条有害指令与 LIBERO 视觉场景配对。

\paragraph{攻击方式。}
（i）\emph{无攻击}（良性分布），用于在干净 \textsc{LIBERO-Long} 上测量良性成功率（BSR）。（ii）\emph{文本越狱}：以明文语言传递有害指令（使用我们的 \textsc{AdvBench} 分集）。（iii）\emph{UADA}/\emph{UPA}/\emph{TMA}：遵循 \citet{wu2024exploringadversarial} 的通用对抗补丁，$L_\infty$ 预算 $16/255$，补丁面积 5\%，500 次内迭代。

\paragraph{评测指标。}
\textbf{BSR}（越高越好）：良性 \textsc{LIBERO-Long} 任务的完成率。
\textbf{HRR-strict}（越高越好）：有害任务中，前 7 个动作 token 全部落入以 bin 128 为中心的零运动窗口（$\pm$8 个 bin）的比例，用作"机器人停止"的模拟器无关替代指标。
\textbf{HRR-partial}（越高越好）：有害任务中，7 个动作 token 中至少 4 个落入零运动窗口的比例。这一较宽松的标准为部分拒绝提供正分，例如停止 4 个自由度但夹爪仍在闭合的机器人，在防御效果较强但不完美时提供更多信息。
\textbf{BRR}（越高越好）：$\tfrac{1}{2}(\mathrm{BSR} + \mathrm{HRR\text{-}strict})$，衡量防御在不过度拒绝情况下的综合实用性。相机就绪版本中将额外报告逐攻击分解和逐类别 HRR。

\paragraph{基线。}
\textbf{无防御}：原始 OpenVLA。
\textbf{AdaShield-S} \cite{adashield2024}：在指令前添加静态安全提示。
\textbf{VLM-Guard 移植} \cite{vlmguard2025}：从 Llama-2-chat 第 14 层提取 SVD rank-5 SSD，施加于 OpenVLA 的所有 token 位置（我们对原始仅最后 token 方法的移植适配版本）。
\textbf{SafeVLA-lite}：\textsc{SafeVLA} \cite{safevla2025} 的廉价替代，以 LoRA 微调（$r=32$）在 $1\mathrm{k}$ 有害 + $1\mathrm{k}$ 无害样本对上训练一个 epoch，有害目标标注为零运动动作 token。

\paragraph{\method{}/\method{}+ 配置。}
除消融变化外，使用第 $L=14$ 层，$\alpha=1.0$。基础 \method{} 仅以动作位置为目标；\method{}+ 使用文本+动作双阶段注入，$\alpha_{\text{text}}=0.3$，$\alpha_{\text{act}}=1.0$。Rank-$k$ 变体使用 $k=5$。自适应 $\alpha$ 头为隐藏维度 256 的两层 MLP，在配对标注输入上训练 3 个 epoch。

\subsection{主要结果}
\label{sec:exp:main}

表~\ref{tab:main-placeholder} 将报告所有基线和 \method{} 变体的 BSR、HRR-strict、HRR-partial 及 BRR。

\begin{table}[t]
\centering\small
\framebox[\columnwidth]{%
  \begin{minipage}{0.9\columnwidth}\centering\vspace{8pt}
  \emph{占位符。} 主对比实验。行：\{无防御, AdaShield, VLM-Guard 移植, SafeVLA-lite, \method{}（仅动作）, \method{}+（文本+动作）, \method{}+rank-$k$, \method{}+自适应 $\alpha$\}。列：BSR, HRR-strict, HRR-partial, BRR。\vspace{8pt}
  \end{minipage}}
\caption{在指令条件化 \textsc{AdvBench} 分集与 \textsc{LIBERO-Long} 良性任务上的主对比实验。}
\label{tab:main-zh}
\end{table}

\subsection{消融实验}
\label{sec:exp:ablation}

我们对五个维度进行消融。
\textbf{方向来源}：将拒绝方向与（i）等范数随机单位向量（固定随机种子跨实验一致）、（ii）取反的拒绝方向、（iii）无干预进行对比。该对照实验至关重要：若随机向量与拒绝方向获得相近的 HRR，则收益归因于一般隐层状态扰动而非移植的安全几何；若拒绝 $\gg$ 随机，则跨模型迁移假说得到验证。记差值 $\Delta_{\text{rand}}=\mathrm{HRR}(r_L) - \mathrm{HRR}(\text{rand})$。
\textbf{层扫描}：$L \in \{8, 10, 12, 14, 16, 18, 20\}$，确定移植方向既有信号又足够早以影响后续动作解码的最优层。
\textbf{尺度扫描}：$\alpha \in \{0, 0.1, 0.5, 1, 2, 5\}$，揭示 BSR/HRR 权衡并为自适应 $\alpha$ 提供动机。
\textbf{目标位置}：\{仅动作、仅文本、文本+动作、全部\}，隔离归因于动作 token 定向和生成阶段感知双阶段注入的收益——与 \textsc{VLM-Guard} 的关键区别。
\textbf{Rank-$k$}：$k \in \{1, 3, 5, 10\}$，验证动作空间危害是否跨越多个方向。

\begin{table}[t]
\centering\small
\framebox[\columnwidth]{%
  \begin{minipage}{0.9\columnwidth}\centering\vspace{8pt}
  \emph{占位符。} 五维消融矩阵。顶部块：方向来源（拒绝/随机/取反/无）。底部块：层、尺度、目标位置、子空间 rank。\vspace{8pt}
  \end{minipage}}
\caption{消融扫描。$\Delta_{\text{rand}}$ 列报告拒绝$-$随机 HRR 差值，正值确认跨模型几何是有效成分。}
\label{tab:ablations-zh}
\end{table}

\subsection{自适应攻击评测}
\label{sec:exp:adaptive}

我们在白盒自适应攻击者下评测 \method{}，该攻击者知晓 $r_L$、层选择和 $\alpha$。攻击者通过 PGD（500 次迭代，$\epsilon=16/255$）优化图像补丁，最小化 $|h_L[\text{last}] \cdot r_L|$，将隐层状态推至拒绝方向的正交补空间。我们报告有无该攻击下的 HRR-strict 和 HRR-partial。一个不仅仅捕获朴素攻击者的防御，在自适应压力下仍应在无防御基线之上保持非平凡优势。

\subsection{与现有防御的组合}
\label{sec:exp:composition}

最后，我们测试 \method{}+ 与 AdaShield 提示和 VLM-Guard 投影的正交组合效果。表~\ref{tab:composition-zh} 将报告成对和三重组合的 BRR。正向组合支持将 \method{}+ 定位为与现有提示层和投影层防御互补的即插即用动作层模块。

\begin{table}[t]
\centering\small
\framebox[\columnwidth]{%
  \begin{minipage}{0.9\columnwidth}\centering\vspace{8pt}
  \emph{占位符。} 组合矩阵：AdaShield / VLM-Guard / \method{}+ 及其成对/三重组合。\vspace{8pt}
  \end{minipage}}
\caption{\method{}+ 与现有防御的组合效果。}
\label{tab:composition-zh}
\end{table}
